\documentclass[conference]{IEEEtran}
\IEEEoverridecommandlockouts

\usepackage{cite}
\usepackage{amsmath}
\usepackage{amssymb}
\usepackage{graphicx}
\usepackage{booktabs}
\usepackage{url}

\def\BibTeX{{\rm B\kern-.05em{\sc i\kern-.025em b}\kern-.08em
    T\kern-.1667em\lower.7ex\hbox{E}\kern-.125emX}}

\begin{document}

\title{Where Does Intent Get Lost? An Offline Trajectory-Matching Analysis of\\
Reasoning-to-Action Handoff Failures in a Humanoid VLA}

% Double-anonymous: no author names/affiliations in the submitted version.
% Fill in the real author block below only for the camera-ready version.
\author{
\IEEEauthorblockN{Anonymous Author(s)}
\IEEEauthorblockA{Submission for double-anonymous review}
}
% \author{
% \IEEEauthorblockN{First Author\IEEEauthorrefmark{1}, Second Author\IEEEauthorrefmark{1}}
% \IEEEauthorblockA{\IEEEauthorrefmark{1}Affiliation, City, Country \\ email@example.com}
% }

\maketitle

\begin{abstract}
Vision-Language-Action (VLA) models for humanoid manipulation typically decompose into a
high-level reasoning stage that plans a sequence of sub-goals and a low-level action stage
that executes continuous joint-space trajectories. Failures in these systems are often
attributed vaguely to ``the model,'' without distinguishing whether the plan itself was
wrong or whether a correct plan failed to survive the handoff into execution. We present
an offline evaluation methodology that scores these two stages independently against real
teleoperated demonstrations and cross-references them into a joint classification,
isolating the specific failure mode of \emph{intent lost in handoff}: a correct plan whose
execution diverges from the demonstrated trajectory. Applied to GR00T N1.7 (zero-shot on a
real Unitree G1 teleop dataset) with Cosmos-Reason2-2B as a reasoning-stage proxy, across
50 pick-and-place episodes (150 sub-goal-phase observations) we find that 68.0\% of phase
observations show intent lost in handoff versus only 3.3\% where both stages fail
together, with failure concentrated in continuous joint control (arm, wrist, or torso
deviates from the tightest tolerance in 100\% of observations) rather than discrete
grasp-state prediction (deviates in only 18.7\%), and rising monotonically through the
task from 58.0\% at the initial reach to 74.0\% by the final retreat. We discuss the
methodological limits of trajectory-matching against a single demonstration as a proxy
for task success, and release the scoring pipeline as an open, auditable tool for
pre-deployment diagnosis of VLA planning-execution failures.
\end{abstract}

\begin{IEEEkeywords}
vision-language-action models, humanoid robots, failure analysis, offline evaluation, imitation learning
\end{IEEEkeywords}

\section{Introduction}

Vision-language-action models increasingly follow a two-stage architecture explicitly
described by their own authors in dual-process terms: a ``System 2'' component performs
deliberate, language-mediated reasoning to decompose a task instruction into sub-goals,
while a ``System 1'' component maps the current observation and sub-goal onto continuous
low-level actions. GR00T N1~\cite{bjorck2025groot} and $\pi_0$~\cite{black2024pi0} both
adopt this framing independently, and Embodied Chain-of-Thought~\cite{zawalski2024ecot}
shows that eliciting an inspectable intermediate reasoning trace, even in an otherwise
single-stage policy, improves both performance and interpretability. This decomposition
is attractive for modularity, but it introduces a specific failure surface that
single-stage, end-to-end policies~\cite{kim2024openvla,brohan2023rt2} do not have by
construction: the reasoning stage can produce a completely correct plan that the action
stage nonetheless fails to execute faithfully. We call this failure mode
\textbf{intent lost in handoff}.

Existing evaluation practice for VLA models typically reports aggregate task success
rate, which conflates planning failures with execution failures and gives no diagnostic
signal about \emph{where} in the reasoning-to-action pipeline a given failure originates.
This is a meaningful gap for humanoid platforms specifically, where whole-body
coordination (balance, torso posture, bimanual coordination) adds execution-stage failure
modes beyond what tabletop single-arm manipulation exposes, and where physical
trial-and-error evaluation is expensive and slow compared to software iteration.

We contribute:

\begin{enumerate}
\item An offline evaluation methodology that scores reasoning-stage and execution-stage
  fidelity independently, against ground truth derived directly from real teleoperated
  demonstrations (no simulation, no physical execution required), and cross-references
  them into a $2\times2$ classification that isolates intent-lost-in-handoff from cases
  where both stages fail together or where an incorrect plan is accidentally executed
  correctly.
\item A concrete instantiation of this methodology on GR00T N1.7 (NVIDIA's whole-body
  humanoid VLA policy) evaluated zero-shot against real Unitree G1 teleoperated
  demonstrations, using Cosmos-Reason2-2B as an explicit, labeled proxy for the reasoning
  stage GR00T does not expose in human-readable form.
\item An empirical failure taxonomy across 50 episodes (150 sub-goal-phase observations)
  showing that execution failures concentrate in continuous joint-space control rather
  than discrete grasp-state prediction, and that failure rate increases through the
  course of a pick-and-place task rather than being uniform.
\item Release of the full scoring pipeline as open, auditable tooling, including the
  specific numeric thresholds used and the reasoning behind design decisions made after
  inspecting real per-episode data (e.g., scoring against an early prediction window
  rather than a full open-loop rollout, to isolate immediate plan-following from
  accumulated open-loop drift).
\end{enumerate}

\section{Related Work}

\textbf{VLA models and dual-system architectures.} GR00T N1~\cite{bjorck2025groot}, the
model we evaluate, explicitly frames itself as a dual-system architecture: a
vision-language module (its authors' ``System 2'') interpreting instructions and scene,
coupled to a diffusion-transformer action module (``System 1'') generating continuous
motor commands. $\pi_0$~\cite{black2024pi0} independently converges on the same framing,
pairing a discrete-token planning stage with a flow-matching action stage. Not every VLA
makes this split explicit: OpenVLA~\cite{kim2024openvla} and
RT-2~\cite{brohan2023rt2} are single-stage policies that map vision and language directly
to action tokens without a separately inspectable planning representation. This
distinction matters for our methodology, which requires a model (or a reasonable proxy
for one) that produces an inspectable plan distinct from its low-level action output; it
is not directly applicable to single-stage VLAs without first eliciting an analogous
reasoning trace some other way.

\textbf{Explicit reasoning and chain-of-thought in embodied models.} Embodied
Chain-of-Thought (ECoT)~\cite{zawalski2024ecot} trains OpenVLA to produce an inspectable
reasoning trace -- plans, sub-tasks, and grounded features like bounding boxes -- before
its action head runs, improving both success rate and interpretability. Our approach is
related but serves a different purpose: ECoT modifies training to make an otherwise
single-stage model produce visible reasoning as part of improving that same model, while
we evaluate an already-dual-system, already-trained model's plan-to-action fidelity
without any additional training, using reasoning and action as independently scored
diagnostic signals rather than as a joint training target.
Cosmos-Reason1~\cite{azzolini2025cosmosreason1} -- the predecessor to Cosmos-Reason2-2B,
which we use as our reasoning-stage proxy (Section~\ref{sec:method}) -- is exactly the
kind of standalone embodied-reasoning VLM this diagnostic approach depends on: a model
that shares architectural lineage with GR00T's internal VLM but exposes natural-language
chain-of-thought output that GR00T's own latent reasoning does not.

\textbf{Failure analysis for robot policies.} RoboFAC~\cite{ye2025robofac} is the closest
prior work in spirit: a large-scale, failure-centric framework (9{,}440 erroneous
trajectories, 78{,}623 QA pairs, simulation and real-world) that trains a dedicated model
to diagnose and correct manipulation failures from trajectory data. It differs from our
approach in two ways relevant to this paper's contribution: it categorizes failures
without explicitly cross-referencing an independently-scored plan against
independently-scored execution the way our $2\times2$ classification does, and it
requires training a diagnostic model on a large labeled failure dataset, whereas our
method is a lightweight, threshold-based offline rubric requiring no additional training
data or model training, specifically designed to isolate the reasoning-to-action handoff
as its own failure surface.

\textbf{Evaluation methodology beyond success rate.} Kress-Gazit et
al.~\cite{kressgazit2024empirical} argue that robot learning research over-relies on
aggregate success rate reported with minimal experimental context, and propose richer
practices: explicit conditions, multiple complementary metrics, statistical rigor, and
qualitative failure description -- demonstrated through physical hardware trials. We
share this motivation but differ in setting: we obtain much of this richer diagnostic
signal entirely offline against existing demonstration data, directly addressing the cost
and slowness of physical trial-and-error evaluation that motivates their call for better
practices in the first place.

\textbf{Humanoid whole-body benchmarks and datasets.}
HumanoidBench~\cite{sferrazza2024humanoidbench} is a simulated benchmark spanning
humanoid locomotion and manipulation tasks; Humanoid Everyday~\cite{zhao2025humanoideveryday}
is a large real-world dataset for diverse humanoid manipulation. Both establish task and
data infrastructure for humanoid learning broadly; neither proposes a
reasoning-versus-execution diagnostic methodology, which is the specific gap this paper
targets. Like Humanoid Everyday, we use real (not simulated) demonstration data, but
apply it to a narrower diagnostic question -- where a specific model's plan-to-action
fidelity breaks down -- rather than broad task or skill coverage.

\section{Method}
\label{sec:method}

\subsection{Dataset}

We use \texttt{nvidia/PhysicalAI-Robotics-GR00T-Teleop-G1}, a public dataset of real (not
simulated) teleoperated Unitree G1 humanoid trajectories in LeRobot format, sampled at
20\,fps. We evaluate on the \texttt{g1-pick-apple} task (pick up a red apple, place it on
a plate). \textit{[TODO: note if additional task folders from this dataset -- pear,
grapes, starfruit -- are added before submission for cross-task generalization.]} State
and action are 43-dimensional whole-body vectors (legs, waist, both arms, both hands).

\subsection{Models}

\textbf{Action stage (System 1):} GR00T N1.7 (3B parameters), a later checkpoint of
NVIDIA's GR00T N1 whole-body humanoid VLA policy~\cite{bjorck2025groot}, loaded zero-shot
with \texttt{EmbodimentTag.REAL\_G1} -- a pretrain tag matched by construction to this
dataset's embodiment, requiring no fine-tuning.

\textbf{Reasoning stage (System 2) proxy:} GR00T's own internal reasoning representation
is latent (hidden-state activations), not human-readable. We use Cosmos-Reason2-2B, a
vision-language model sharing backbone lineage with GR00T's internal VLM but released
standalone with natural-language chain-of-thought output, run independently on the same
task instruction and initial frame. Cosmos-Reason2 is the successor to
Cosmos-Reason1~\cite{azzolini2025cosmosreason1} and, as of this writing, is documented via
model card and code release rather than a dedicated paper; we cite the Reason1 paper for
the underlying architecture and training approach the Reason2 line descends from.
\textbf{This is an explicit proxy, not GR00T's literal internal state}, and all results
involving the reasoning stage must be read with this caveat.

\subsection{Ground-truth sub-goal segmentation}

The dataset provides one flat task instruction per episode, not labeled sub-goal
boundaries. We derive reach / transport / retreat sub-goal boundaries from the recorded
hand-joint trajectory: a scalar ``hand closure'' signal (mean absolute value of the 7
hand-joint state dimensions per frame, per hand) is computed for both hands; the hand
with the larger observed range is taken as the grasping hand for that episode (this is
computed from ground-truth recorded data, independent of any model prediction). A grasp
event is the first sustained ($\geq 3$-frame) crossing of a per-episode-fit threshold
from below to above; a release event is the following sustained reverse crossing. This
splits each episode into reach (start $\rightarrow$ grasp), transport (grasp
$\rightarrow$ release), and retreat (release $\rightarrow$ end).

\subsection{Execution-stage scoring}

For each of the three sub-goal phases, GR00T is queried once at that phase's start frame,
producing a predicted action chunk. This is compared against the actual recorded
teleoperated action for the same window, per body-part group: wrist position (Euclidean
distance, cm), wrist rotation (geodesic angle between predicted and actual rotation
matrices, reconstructed from the 6D continuity representation of Zhou et al.), arm/torso
joint angles (per-joint degree difference), and hand grasp state (discrete
open/transitioning/closed bucket match). Table~\ref{tab:thresholds} gives the
match/minor-deviation/failure thresholds per dimension.

\begin{table}[htbp]
\centering
\caption{Execution-stage scoring thresholds}
\label{tab:thresholds}
\begin{tabular}{@{}lccc@{}}
\toprule
Dimension & Match & Minor deviation & Failure \\
\midrule
Wrist position & $<3$\,cm & 3--7\,cm & $>7$\,cm \\
Wrist rotation & $<10^\circ$ & $10$--$20^\circ$ & $>20^\circ$ \\
Arm/waist joint angle & $<5^\circ$ & $5$--$15^\circ$ & $>15^\circ$ \\
\bottomrule
\end{tabular}
\end{table}

A body-part group (e.g., ``left arm,'' 7 joints) is scored \emph{match} if all joints are
match-tier, \emph{failure} if any single joint is failure-tier, \emph{minor deviation}
otherwise. Only the arm, hand, wrist, and torso of the \emph{grasping} side (plus the
shared torso/waist) drive the phase-level verdict; the idle arm's motion is not commanded
by the task and its prediction error is not evidence of a failed handoff.

Scoring uses only the first 10 predicted steps ($\sim$0.5\,s at this dataset's frame
rate), not the full $\sim$40-step predicted chunk. Inspecting signed per-joint error on
real data showed two distinct phenomena conflated by full-chunk averaging: some joints
show error that grows steadily from near-zero over the open-loop horizon (expected
behavior for any policy evaluated this way, since deployment normally re-plans every few
steps rather than executing a long open-loop rollout), while others show a near-constant
offset present even at the first predicted step (a more direct signal about whether the
immediate action reflects the plan). We report both windows; the early window drives
classification.

\subsection{Reasoning-stage scoring}

Cosmos-Reason2's step-by-step plan for the task is evaluated once per episode (it
produces one plan up front, not a new plan per phase) via LLM-as-judge: does the plan
cover all three canonical sub-goal phases, in the correct order, naming the correct
object and target, without hallucinated content or being too vague to map onto the
phases. We observed some inconsistency in the judge's handling of borderline elaborations
across calls (Section~\ref{sec:discussion}); this is a known limitation of LLM-as-judge
approaches generally.

\subsection{Joint classification}

Per phase, we cross-reference reasoning-stage match against execution-stage tier
(Table~\ref{tab:classification}). ``Minor deviation'' execution results are tracked as a
separate category rather than folded into match or failure.

\begin{table}[htbp]
\centering
\caption{The core reasoning $\times$ execution classification}
\label{tab:classification}
\begin{tabular}{@{}lcc@{}}
\toprule
 & Execution match & Execution failure \\
\midrule
Reasoning match & Success & \textbf{Intent lost in handoff} \\
Reasoning failure & Accidentally correct & Compounding failure \\
\bottomrule
\end{tabular}
\end{table}

\section{Results}

Across 50 episodes (150 sub-goal-phase observations, all with valid grasp/release
segmentation), Table~\ref{tab:overall} gives the overall classification breakdown and
Table~\ref{tab:byphase} the breakdown by sub-goal phase.

\begin{table}[htbp]
\centering
\caption{Overall classification}
\label{tab:overall}
\begin{tabular}{@{}lrr@{}}
\toprule
Classification & Count & \% \\
\midrule
Intent lost in handoff & 102 & 68.0\% \\
Minor deviation & 43 & 28.7\% \\
Compounding failure & 5 & 3.3\% \\
Success & 0 & 0.0\% \\
Accidentally correct & 0 & 0.0\% \\
\bottomrule
\end{tabular}
\end{table}

\begin{table}[htbp]
\centering
\caption{Classification by sub-goal phase}
\label{tab:byphase}
\begin{tabular}{@{}lrrr@{}}
\toprule
Phase & Intent lost & Minor dev. & Compounding \\
\midrule
Reach & 58.0\% & 42.0\% & 0.0\% \\
Transport & 72.0\% & 26.0\% & 2.0\% \\
Retreat & 74.0\% & 18.0\% & 8.0\% \\
\bottomrule
\end{tabular}
\end{table}

Failure rate rises monotonically through the task: the model's execution is closest to
the demonstrated trajectory immediately after the observation (reach), and diverges
further as the task progresses through transport into retreat.

Table~\ref{tab:driving} shows which measured dimension most often drives a
failure/minor-deviation verdict (a phase can implicate more than one dimension).
Discrete grasp-state prediction (hand open/closed) is reliable; continuous joint-space
and pose control (arm, wrist, torso) is where deviation concentrates. This suggests the
model's difficulty is specifically in precise continuous control fidelity, not in
higher-level object-interaction state tracking.

\begin{table}[htbp]
\centering
\caption{Which measured dimension drives failure/minor-deviation (of 150)}
\label{tab:driving}
\begin{tabular}{@{}lr@{}}
\toprule
Dimension & Count \\
\midrule
Waist (torso) rotation & 148 \\
Arm joint angles & 145 \\
Wrist position & 145 \\
Wrist rotation & 100 \\
Hand grasp state & 28 \\
\bottomrule
\end{tabular}
\end{table}

Table~\ref{tab:reasoning} gives reasoning-stage failure categories at the episode level.
The dominant failure pattern is a correct plan with a failed execution
(\texttt{intent\_lost\_in\_handoff}, 68\% of observations), not a failure of both stages
together (\texttt{compounding\_failure}, 3.3\%) -- consistent with the reasoning stage
(or its proxy) being comparatively reliable relative to execution-stage trajectory
fidelity in this setting. We manually verified all 3 \texttt{compounding\_failure}
episodes against their raw reasoning traces and execution numbers: in two, the plan never
described a retreat/withdraw step at all; in the third, the plan described the
\emph{plate} moving to meet the apple rather than the reverse -- a clear reversal of
actor and object, not a borderline judge call.

\begin{table}[htbp]
\centering
\caption{Reasoning-stage failure categories (episode-level, $n=50$)}
\label{tab:reasoning}
\begin{tabular}{@{}lr@{}}
\toprule
Category & Episodes \\
\midrule
Hallucinated step & 10 (20\%) \\
Wrong object/target & 9 (18\%) \\
Missing sub-goal & 5 (10\%) \\
\bottomrule
\end{tabular}
\end{table}

Table~\ref{tab:matchcount} shows how many of the 5 scored dimensions (arm, hand, wrist
position, wrist rotation, waist) land in the strictest match tier simultaneously, per
phase. No phase in our sample achieves the strictest tier on all 5 dimensions
simultaneously, which explains the 0\% strict-success rate in Table~\ref{tab:overall} --
but the shape of this distribution is itself informative. If a single overly strict
threshold were responsible (e.g., the torso-rotation tolerance being unreasonably tight
relative to what a well-performing model could achieve), we would expect phases to
cluster near 4/5, with one holdout dimension. Instead the distribution is centered at
1--2/5 (82\% of phases), with only 4 of 150 phases (2.7\%) reaching 3 or more. This is
consistent with genuinely independent, non-trivial deviation across multiple
continuous-control dimensions at once, not an artifact of one miscalibrated cutoff.

\begin{table}[htbp]
\centering
\caption{Dimensions matching simultaneously, per phase (of 5)}
\label{tab:matchcount}
\begin{tabular}{@{}lrr@{}}
\toprule
Dimensions matching & Count & \% \\
\midrule
0 of 5 & 23 & 15.3\% \\
1 of 5 & 75 & 50.0\% \\
2 of 5 & 48 & 32.0\% \\
3 of 5 & 3 & 2.0\% \\
4 of 5 & 1 & 0.7\% \\
5 of 5 & 0 & 0.0\% \\
\bottomrule
\end{tabular}
\end{table}

% [TODO Task 3: figures -- bar chart of Table~\ref{tab:overall}, stacked bar of
% Table~\ref{tab:byphase} by phase, bar chart of Table~\ref{tab:matchcount}, example
% frame sequence for a representative intent-lost-in-handoff case with
% predicted-vs-actual overlay if time permits]

\section{Discussion and Limitations}
\label{sec:discussion}

\textbf{Trajectory-matching is not task-success measurement.} This is the central
methodological caveat of the entire approach and must not be elided: every
classification here measures divergence from \emph{one recorded human demonstration},
not whether the predicted trajectory would have physically succeeded at the task. A
pick-and-place task generally admits multiple valid solution paths; GR00T predicting a
different-but-valid grasp approach would be scored identically to a genuinely failed
grasp under this rubric. The zero success rate we observe (Table~\ref{tab:overall})
should be read as ``GR00T's zero-shot predictions on this embodiment do not closely track
this specific demonstrated trajectory,'' not as ``GR00T fails to complete this task.''
Table~\ref{tab:matchcount}'s distribution supports this being a real measurement of
partial trajectory fidelity rather than a rubric artifact -- most phases achieve 1--2 of
5 dimensions within the strictest tolerance, not zero, and the shape of the distribution
is inconsistent with one miscalibrated threshold being solely responsible. Grounding the
trajectory-fidelity-vs-task-success distinction -- ideally via the physical-consequence
proxy layer described below -- is the most important piece of future work this result
motivates.

\textbf{Reasoning-stage proxy.} Cosmos-Reason2-2B is used as an explicit, labeled stand-in
for GR00T's internal reasoning, which is otherwise a latent, non-human-readable
representation. Results describing ``reasoning-stage'' behavior describe this proxy, not
a literal decoding of GR00T's internals.

\textbf{Forward-kinematics approximation.} GR00T's \texttt{REAL\_G1} embodiment expects
wrist end-effector pose (\texttt{wrist\_eef\_9d}) state/action fields that no public G1
dataset, including this one, records directly. We compute these via forward kinematics
from the official Unitree G1 URDF. This is an approximation not independently verified
against NVIDIA's internal training convention for this field; any bias in that convention
should largely cancel out of the \emph{relative} predicted-vs-actual comparisons this
rubric performs (both sides are computed the same way), but this has not been
independently confirmed.

\textbf{LLM-as-judge noise.} The reasoning-stage judge showed some inconsistency across
calls in borderline cases -- e.g., an unspecified color descriptor (``teal plate'') was
judged harmless elaboration in one episode and flagged as hallucination in others. This
is a known limitation of LLM-as-judge approaches generally rather than specific to this
pipeline, and a source of noise in the reasoning-failure-category counts
(Table~\ref{tab:reasoning}) specifically, though it does not affect the execution-stage
numbers (computed deterministically from numeric thresholds). We manually re-verified all
3 \texttt{compounding\_failure} episodes -- the cases most sensitive to judge reliability,
since both stages must be independently correct for that label to be trustworthy --
against the raw reasoning trace and execution numbers directly; all 3 held up as genuine
double-failures rather than judge noise.

\textbf{Single task, single object.} Results here are from one task type
(\texttt{g1-pick-apple}). \textit{[TODO: update if additional task folders are added
before submission.]} The dataset includes three other object types (pear, grapes,
starfruit) under the same task structure; whether the phase-wise failure gradient and the
discrete/continuous split generalize across object geometry is untested.

\textbf{Sample size.} 50 episodes / 150 phase observations is a modest sample for a
strong generalization claim; error bars and statistical testing \textit{[TODO if time
permits]} would strengthen Tables~\ref{tab:overall}--\ref{tab:reasoning}'s claims, in the
spirit of the statistical-rigor practices Kress-Gazit et al.~\cite{kressgazit2024empirical}
argue robot learning evaluation should adopt more broadly.

\textbf{Physical-consequence proxy layer (explicitly out of scope here).} A natural
extension is a kinematic-plausibility layer estimating whether a flagged execution
failure plausibly causes a physical consequence (e.g., a balance-relevant torso
deviation, an unreachable target) -- applied on top of, not in place of, the
match/failure classification above. We deliberately scoped this out of the current
submission to prioritize evaluating at a larger episode count within the available time;
it remains the most direct way to address the trajectory-matching-vs-task-success gap
above and is future work.

\section{Conclusion}

\textit{[TODO: 1 paragraph restating the contribution and the headline empirical finding
once final numbers are locked in.]}

\bibliographystyle{IEEEtran}
\bibliography{references}

\end{document}
